% ---------------------------------------------------
% Trabajo Final de Grado
% Author: Fabián González Lence <alu0101549491@ull.edu.es>
% Chapter: Metodología
% ----------------------------------------------------

\chapter{Metodología} \label{chap:Metodology}

El presente capítulo describe en detalle la metodología adoptada para el desarrollo experimental del trabajo.
Se articulan los cinco proyectos que conforman el experimento, el flujo de trabajo multimodelo con sus cuatro fases, las plantillas de prompts empleadas, las métricas de evaluación definidas y la planificación temporal del proceso.\\

\section{Diseño experimental} \label{sec:diseno-experimental}

El experimento que vertebra este trabajo se articula en torno a cinco proyectos web de complejidad creciente, cada uno de los cuales actúa como caso de estudio independiente dentro de un mismo marco metodológico.
La progresión de complejidad no es arbitraria: cada proyecto introduce uno o varios elementos nuevos ---arquitectónicos, tecnológicos o de escala--- que no estaban presentes en el anterior, permitiendo observar cómo las capacidades y limitaciones de los modelos de \ac{IA} se manifiestan de forma diferencial según el nivel de exigencia del sistema a desarrollar.\\

La Tabla~\ref{tab:proyectos-resumen} recoge las características principales de los cinco proyectos, incluyendo la arquitectura adoptada, la pila tecnológica, el número aproximado de ficheros de producción y el tipo de pruebas generadas.\\

\begin{table}[H]
\centering
\small
\renewcommand{\arraystretch}{1.3}

\begin{tabularx}{\textwidth}{|p{1.5cm}|p{2.5cm}|Y|Y|p{1.5cm}|p{2cm}|}
\hline
\textbf{Proyecto} & \textbf{Arquitectura} & \textbf{Tecnologías} & \textbf{Función clave} & \textbf{Ficheros aprox.} & \textbf{Testing} \\
\hline

\ac{HANGMAN} &
\ac{MVC} + Composite, Observer &
TypeScript, Vite, Canvas \ac{API}, Bulma, Jest &
Juego de letras con renderizado en Canvas, diccionario externo &
9 módulos &
Jest unitario (AAA) \\
\hline

\ac{PLAYER} &
Component-Based + Custom Hooks &
React 18, TypeScript, Vite, Jest, React Testing Library &
Reproducción de audio, gestión de playlist, persistencia en localStorage &
$\sim$15 módulos &
Jest + RTL unitario \\
\hline

\ac{BALATRO} &
Layered Architecture (MVC + Services) &
React, TypeScript, Vite, Jest &
Manos de póker, progresión por niveles, tienda, cartas especiales, assets generados con IA &
$\sim$30 módulos &
Jest unitario \\
\hline

\ac{CARTO} &
Clean Architecture (4 capas) &
Vue.js 3, TypeScript, Pinia, Socket.io, Axios, Playwright &
Gestión de proyectos, tareas, mensajería en tiempo real, integración Dropbox, roles diferenciados &
$>$300 ficheros &
Playwright E2E \\
\hline

\ac{TENNIS} &
Hexagonal Architecture (frontend + backend separados) &
Angular 19, TypeScript, Node.js, Socket.io, Playwright &
Gestión de torneos, brackets, order of play, sistema de dobles, notificaciones, múltiples roles &
$>$400 ficheros &
Playwright E2E \\
\hline

\end{tabularx}

\caption{Resumen comparativo de los cinco proyectos del experimento.}
\label{tab:proyectos-resumen}
\end{table}

\subsection{Proyecto 1 — The Hangman Game (\ac{HANGMAN})} \label{subsec:p1-hangman}

\ac{HANGMAN} es una \ac{SPA} que implementa el clásico juego del ahorcado en TypeScript puro, sin ningún framework de interfaz.
El jugador debe adivinar palabras de un diccionario de animales antes de completar el dibujo del ahorcado en un canvas HTML5, con un máximo de seis intentos fallidos.
Su arquitectura sigue el patrón \ac{MVC} con aplicación adicional del patrón Composite en la capa de vista ---que compone cuatro componentes especializados: \texttt{WordDisplay}, \texttt{AlphabetDisplay}, \texttt{HangmanRenderer} y \texttt{MessageDisplay}--- y del patrón Observer para la gestión de eventos de usuario.\\

Este proyecto fue concebido como el punto de partida del experimento por varias razones: la ausencia de framework elimina variables externas, el número reducido de módulos (9) lo hace manejable como primer caso de prueba, y la arquitectura \ac{MVC} es suficientemente conocida como para disponer de criterios de evaluación bien establecidos.
El stack tecnológico se limita a TypeScript, Vite, la \ac{Canvas API} nativa y Bulma como framework CSS, con Jest y ts-jest para las pruebas unitarias.\\

\subsection{Proyecto 2 — Music Web Player (\ac{PLAYER})} \label{subsec:p2-player}

\ac{PLAYER} es un reproductor de música web construido con React 18 y TypeScript que permite la reproducción de canciones locales, la gestión de una lista de reproducción y la visualización de información de cada pista.
La arquitectura adopta el patrón de Component-Based Architecture con Custom Hooks, donde la lógica de estado se encapsula en tres hooks reutilizables: \texttt{useAudioPlayer} (integración con la \ac{API} de Audio HTML5), \texttt{usePlaylist} (gestión de la lista de reproducción) y \texttt{useLocalStorage} (persistencia entre sesiones).\\

Respecto a \ac{HANGMAN}, este proyecto introduce tres elementos nuevos: el uso de React como framework declarativo de interfaz, la separación entre componentes contenedor y componentes presentacionales, y la necesidad de gestionar estado asíncrono asociado a eventos multimedia.
El framework de testing combina Jest con React Testing Library, lo que exige pruebas orientadas al comportamiento visible del usuario en lugar de a la implementación interna.\\

\subsection{Proyecto 3 — Mini Balatro (\ac{BALATRO})} \label{subsec:p3-balatro}

\ac{BALATRO} es un juego de cartas inspirado en el videojuego Balatro, en el que el jugador debe superar niveles de puntuación jugando manos de póker con una baraja francesa de 52 cartas, complementadas por cartas especiales (planetas, tarot y jokers).
El sistema incluye progresión por niveles, encuentros con jefes cada tres rondas, economía de compra en tienda y persistencia del estado de partida en localStorage.\\

La arquitectura adopta una Layered Architecture que combina el patrón \ac{MVC} con una capa de servicios transversales para la tienda, la configuración del juego y la persistencia.
Este proyecto supone un salto de complejidad lógica significativo respecto a los anteriores: la evaluación de manos de póker, el cálculo del sistema de puntuación chips/multiplicador y la gestión de las interacciones entre cartas especiales requieren un dominio del modelo con muchas más entidades y relaciones que los proyectos previos.
Adicionalmente, los assets gráficos de las cartas fueron generados con Google Gemini, introduciendo por primera vez la dimensión de generación de recursos visuales mediante \ac{IA}.\\

\subsection{Proyecto 4 — Cartographic Project Manager (\ac{CARTO})} \label{subsec:p4-carto}

\ac{CARTO} es una aplicación web y móvil para la gestión de proyectos cartográficos entre un administrador y múltiples clientes simultáneos.
Sus funcionalidades incluyen gestión de proyectos con seguimiento de estado, asignación bidireccional de tareas con cinco estados posibles, mensajería interna por proyecto con adjuntos, integración con Dropbox para ficheros técnicos, sistema de roles con permisos diferenciados y notificaciones en tiempo real mediante WebSockets.\\

La arquitectura adopta los principios de Clean Architecture, organizando el código en cuatro capas: Dominio (entidades, objetos de valor, interfaces de repositorio), Aplicación (servicios e interfaces), Infraestructura (implementaciones de repositorios, cliente HTTP Axios, manejador de Socket.io, servicios externos) y Presentación (componentes Vue.js 3, stores Pinia, composables y router).
Con más de trescientos ficheros distribuidos en estas capas, \ac{CARTO} representa el punto de inflexión en el que el flujo conversacional de los proyectos anteriores resultó inviable, motivando el cambio metodológico descrito en el Capítulo~\ref{chap:AIModels}.
El testing abandona las pruebas unitarias con Jest para adoptar pruebas \textit{end-to-end} con Playwright, más adecuadas para verificar el comportamiento integrado de una aplicación full-stack.\\

\subsection{Proyecto 5 — Tennis Tournament (\ac{TENNIS})} \label{subsec:p5-tennis}

\ac{TENNIS} es una aplicación web responsive para la gestión integral de torneos de tenis, orientada a tres tipos de usuario con roles y permisos diferenciados: organizadores, participantes registrados y público general.
El sistema cubre el ciclo completo de un torneo: creación y configuración, registro de participantes (incluyendo un sistema de invitaciones para modalidad dobles), diseño del cuadro, generación del order of play con gestión de pistas y horarios, registro de resultados, cálculo de clasificaciones y publicación de comunicados.\\

Es el proyecto arquitectónicamente más complejo del experimento, con una separación explícita entre frontend y backend como subproyectos dentro del mismo monorepo.
El frontend utiliza Angular 19 con componentes standalone, mientras que el backend está implementado en Node.js sobre Express e incorpora Socket.io para la comunicación en tiempo real.
La arquitectura del frontend sigue una organización inspirada en principios hexagonales, con capas de dominio, aplicación, infraestructura y presentación.
El agente de codificación planificó 24 categorías de ficheros y generó en torno a 420 ficheros de código fuente de producción, lo que lo convierte en el proyecto de mayor escala del trabajo.\\

\subsection{Justificación de la progresión de complejidad} \label{subsec:justificacion-progresion}

La secuencia de los cinco proyectos fue diseñada para que cada uno introdujera un incremento de complejidad manejable pero significativo respecto al anterior, de forma que las limitaciones de los modelos se manifestasen de forma gradual y observable.
\ac{HANGMAN} establece la línea base con arquitectura \ac{MVC} plana y TypeScript puro; \ac{PLAYER} introduce el paradigma declarativo de React y la gestión de estado asíncrono; \ac{BALATRO} añade complejidad lógica del dominio y generación de assets visuales; \ac{CARTO} escala a una aplicación full-stack multicapa con tiempo real e integración de servicios externos; y \ac{TENNIS} maximiza la complejidad con una arquitectura desacoplada frontend/backend, Angular 19 y funcionalidades de negocio sofisticadas.\\

Esta progresión responde también a una lógica de aprendizaje metodológico: cada proyecto permitió refinar los prompts, ajustar los criterios de evaluación y detectar los patrones de fallo de los modelos antes de enfrentarse al siguiente nivel de exigencia.
El cambio de herramientas entre los proyectos 3 y 4 ---de Mistral/Qwen a GitHub Copilot agente--- no fue planificado desde el inicio, sino que surgió precisamente como respuesta a las limitaciones observadas durante el desarrollo de \ac{BALATRO} y las primeras fases de \ac{CARTO}, convirtiéndose en uno de los hallazgos metodológicos más relevantes del trabajo.\\